\section{Related work}
\label{sec.related}

\paragraph{Common prefix}

The backbone line of work~\cite{gkl15,gkl24} abstracts Bitcoin as a protocol
with three properties and proves them under honest majority. The one this paper
revisits is common prefix, which holds with a parameter $k$ of the order of
$\lambda f$~\cite[Theorem 4.10]{gkl24}. The parameter quantifies over pairs of
honest parties, and both of them mine and are synchronised. The same
restriction is present in the bounded delay analysis of~\cite{pss16}, in the
speed and security tradeoffs of~\cite{kp15}, and in the composable treatment
of~\cite{bmtz17}. We import the model of~\cite{gkl24} and its typical execution
machinery unchanged, and we replace the quantifier: one of the two chains is
held by a party that does not mine and is not synchronised with anyone.

\paragraph{Settlement bounds}

A separate line computes the constant $k$ a well connected observer needs.
\cite{dembo20} reduces every attack on a longest chain protocol to a race and
gives the exact tolerable adversarial fraction as a function of the block
interval and the delay. \cite{grr22,grr23} give concrete and computable
settlement bounds for proof of work and for longest chain consensus, and
\cite{guoren22} gives closed form upper and lower bounds on the probability of
a safety violation as a function of $k$. These results sharpen the constant. We
show that for a client whose own link is unconstrained there is no constant to
sharpen. The two statements do not conflict. Their observer receives every
honest message within the delay bound, and ours receives what the adversary
chooses.

\paragraph{Simplified payment verification}

Section 8 of~\cite{nakamoto08} gives the first client protocol: keep the
headers of the longest chain and a Merkle branch. Its network assumption is the
clause quoted in \Cref{sec.intro}, which asks the client to convince itself
that it holds the longest chain rather than to wait a stated time.
\cite{sz16} revisits Bitcoin's security for different classes of defender and
observes that a client that does not relay is strictly weaker than a full node.
Our construction replaces the clause with a rule the client can evaluate, using
the round and the chain it holds.

\paragraph{Succinct clients}

Superblock proofs~\cite{nipopow}, their extension to
miners~\cite{logspace}, probabilistic sampling~\cite{flyclient}, the
interactive protocol of~\cite{blink} and the bisection game
of~\cite{lazyclients} all shorten what a client downloads, from linear in the
length of the chain to polylogarithmic or constant. The question there is how
to compress the evidence. The question here is how deep to prune once the
evidence is in hand. The two compose. A succinct client that has established
the longest chain among those it was shown must still decide how much of it to
keep. Our client is not succinct.

\paragraph{Bootstrapping}

Ouroboros Genesis~\cite{genesis} is the closest structural analogue we know.
Its chain selection rule drops the moving checkpoint rule, under which a party
rejects a fork more than $k$ blocks deep. In its place it compares the density
of the candidate chains in the $s$ slots that follow their fork point. The
reason is the one that drives this paper. A party joining from genesis, or
returning after a long absence, cannot use a constant $k$. The rules differ in
what they measure. Genesis compares two chains locally after their fork point,
in a proof of stake protocol with a slot clock. We compare one chain against
the round, in proof of work.

\paragraph{Sleepy parties}

The sleepy model~\cite{sleepy} lets honest parties go offline and return, and
proves security under honest majority among the parties that are awake. A party
that wakes receives every message it would have received while asleep, so after
waking it is in the position of a synchronised party again, and a constant
depth serves it. The same is true of the dynamically available protocols built
on that model~\cite{ebbflow,goldfish} and of the accountability gadgets
of~\cite{accountability}. Our client never recovers that position: its link
stays adversarial for the whole execution.

\paragraph{Client models}

\cite{cuamdr} makes the client model a first class parameter and crosses four
binary choices, among them whether a client is sleepy and whether it can send
messages as well as receive them. Longest chain consensus sits in their sleepy
silent client model, where they prove a resilience of just under one half
optimal. Their clients share the network of the validators, and degraded
connectivity is captured by sleepiness, whose cost is that a sleeping client
cannot record when a message arrived. Our client is the complement. It is
silent in their sense, it keeps its clock, and what it gives up is any bound on
its own link. Their Appendix F shows that under adversary majority a client
that confirms a $k$-deep block and gossips it is unsafe for every $k$, by a
balancing attack. \Cref{cor.noconstant} is a different statement: it assumes
honest majority, and it says how fast the depth must grow.
\cite{bsts} recovers safety for sleepy clients during adversary majority, for
protocols that issue certificates. \cite{stubborn} studies a client that
gossips a block before confirming it. Neither route is open here, because
certificates are unavailable in the proof of work setting~\cite{lpr21}.

\paragraph{On-chain clients}

A Bitcoin script cannot verify a chain directly, so bridges build a client out
of a dispute game. The client of~\cite{bitvm2} works in the static difficulty setting. An operator
commits to a recent block height. The recency of the claim is enforced by an
absolute timelock, and the operator must then exhibit headers up to that
height. That is our rule executed on chain. An untrusted party
claims a height, and the claim is checked against a clock the party cannot
move. We check the height a peer claims against the round the client reads off
its own clock. The settings differ in what enforces the check. Theirs is an
economic game with a counterparty to slash, ours is a client with no
counterparty. \cite{bitvm3} shows that the static construction fails under variable
difficulty. It repairs it by binding the committed timestamps to the consensus
time of the chain. Its output is a block that after $u$ rounds is at depth at
least $k$ in every honest chain, which is the guarantee \Cref{cor.prune} gives.
\cite{clementine} treats variable difficulty with a permissioned set of
verifiers, and~\cite{glimpse} gives an on-demand client with constant-size
storage.

\paragraph{Partial synchrony}

Under partial synchrony~\cite{dls88} messages are delivered within a known
bound after an unknown global stabilisation time. The client's link here is
weaker: no bound applies at any round. Safety does not use one, by
\Cref{cor.prune} and \Cref{thm.consistency}. Liveness does, and what a
stabilisation time would buy is a bound on the staleness $\stale$ of
\Cref{def.stale}, which is the only quantity \Cref{thm.liveness} charges for
the network.
